Una espira triangular de $l = 4,0\ \mathrm{m}$ de costat com la de la figura es troba inicialment ($t = 0,0$) situada a una distància de 6,0~m d'una regió on hi ha un camp magnètic $B$ perpendicular al pla del paper i cap endins.

\figura[0.6]{fig1.png}

\begin{apartat}{1}
Indiqueu l'expressió de la FEM induïda a l'espira quan aquesta s'endinsa a la regió on hi ha el camp magnètic. Determineu el valor de $B$ sabent que, per a $t = 4,0\ \mathrm{s}$, la FEM induïda és $E = 160\ \mathrm{V}$.
\end{apartat}

\begin{apartat}{1}
Representeu gràficament la FEM induïda $E = E(\mathrm{t})$ entre $t = 0,0$ i $t = 8,0\ \mathrm{s}$. Indiqueu en cada instant el sentit del corrent induït a l'espira.
\end{apartat}
